\Needspace{8\baselineskip}\originalchapter{2006 作业一: 参数与非线性屏蔽}
题源: Ron Parker, 12 Sep 2006, due 19 Sep. 网站另列2002和2005历史答案, 不是本套当期答案; 本章以下均为独立补解.
\begin{exercise}\label{ex:06-1-1}
原题1. 三维Maxwell $f=n(\sqrt\pi v_t)^{-3}\e^{-v^2/v_t^2}$, $v_t^2=2k_BT/m$. 求 (i)各分量均值; (ii)平均速率; (iii)RMS速率; (iv)单方向动能均值; (v)单向穿过平面通量.
\end{exercise}
\begin{solution}
沿作业\ref{ex:03-1-1}的Gaussian积分, 得 $\langle v_i\rangle=0$, $\langle v\rangle=2v_t/\sqrt\pi$, $\sqrt{\langle v^2\rangle}=\sqrt{3/2}v_t$, $\langle mv_i^2/2\rangle=k_BT/2$, $\Gamma=nv_t/(2\sqrt\pi)$. 此处题T为K, 后文能量T替代kBT时会明确.
\end{solution}
\begin{exercise}\label{ex:06-1-2}
原题2. 以A.Bers课程参数表, 给定下表的log10 n(m$^{-3}$)与log10 Te(K), 补Debye长、\ensuremath{\Lambda}=n\ensuremath{\lambda}D\ensuremath{{}^{3}}、fpe(Hz), 判断强耦合与需量子处理者. 原PDF第3页共有十六类，参数见结果表前两列；高密度、极高温行的经典公式只作形式外推。
\end{exercise}
\begin{solution}
用 $\lambda_D=\sqrt{\epsilon_0k_BT_e/(ne^2)}$, $f_{pe}=\sqrt{ne^2/(\epsilon_0m_e)}/(2\pi)$. 表中\ensuremath{\Lambda}无4\ensuremath{\pi}/3, 不同于Debye球实际粒子数ND.
\begin{center}\small\begin{tabular}{lrrrrr}
\toprule 类别&log n&log TK&\ensuremath{\lambda}D(m)&\ensuremath{\Lambda}&fpe(Hz)\\\midrule
电离层D&9&2.5&.0388&$5.84\,10^4$&$2.84\,10^5$\\
弱电流放电&17&4&$2.18\,10^{-5}$&$1.04\,10^3$&$2.84\,10^9$\\
强电流放电&21&5&$6.90\,10^{-7}$&329&$2.84\,10^{11}$\\
MHD转换器&22&3&$2.18\,10^{-8}$&.104&$8.98\,10^{11}$\\
星际气体&0&3.5&3881&$5.84\,10^{10}$&8.98\\
太阳风&6.5&5&12.27&$5.84\,10^9$&$1.60\,10^4$\\
电离层F2&11.5&3&.00388&$1.85\,10^4$&$5.05\,10^6$\\
日冕&13&6.5&.0388&$5.84\,10^8$&$2.84\,10^7$\\
tokamak&20&8&$6.90\,10^{-5}$&$3.29\,10^7$&$8.98\,10^{10}$\\
碱金属表面电离&18&3&$2.18\,10^{-6}$&10.39&$8.98\,10^9$\\
激光等离子体&25&5&$6.90\,10^{-9}$&3.29&$2.84\,10^{13}$\\
核爆等离子体&26&6&$6.90\,10^{-9}$&32.9&$8.98\,10^{13}$\\
脉冲星磁层&18&16&6.90&$3.29\,10^{20}$&$8.98\,10^9$\\
金属&29&2.5&$3.88\,10^{-12}$&$5.84\,10^{-6}$&$2.84\,10^{15}$\\
恒星内部&33&7.5&$1.23\,10^{-11}$&1.85&$2.84\,10^{17}$\\
白矮星&38&7&$2.18\,10^{-14}$&.00104&$8.98\,10^{19}$\\\bottomrule
\end{tabular}\end{center}
以Wigner半径 $a_W=(3/(4\pi n))^{1/3}$ 定义形式上的经典相关参数$\Gamma=e^2/(4\pi\epsilon_0a_Wk_BT)$，以热德布罗意波长$\lambda_{\rm th}=h/\sqrt{2\pi m_ek_BT}$定义简并指标$\eta=n\lambda_{\rm th}^3$。MHD转换器$\Lambda=.104,\Gamma\simeq.58$，是集体Debye近似失效的中等相关例子；激光等离子体$\Lambda\simeq3.3$、恒星内部$\Lambda\simeq1.85$也不能视为强弱耦合渐近区。金属和白矮星的形式$\Lambda$分别为$5.84\,10^{-6}$和$.00104$，但$\eta$分别约$7.36\,10^3$、$1.31\,10^6$，经典Boltzmann--Debye模型本身失效，必须改用量子简并响应；白矮星还涉及相对论性电子。恒星内部$\eta\simeq2.33$也需量子修正。脉冲星行的$T=10^{16}$ K使非相对论模型不适用，真实磁层还涉及强磁场与对产生。表内这些行的$\lambda_D,\Lambda,f_{pe}$只是把指定$n,T$代入所问经典式所得，不是可靠的天体预测；近距散射的量子截止与统计简并是不同判断。
\end{solution}
\begin{exercise}\label{ex:06-1-3}
原题3. 无穷电荷薄片面密度\ensuremath{\sigma}位于x0, 电子/单离子Boltzmann且Te=Ti=T, 放弃e\ensuremath{\phi}\ensuremath{\ll}T近似. (i)Poisson乘\ensuremath{\phi}\ensuremath{{}^{\prime}}求\ensuremath{\phi}\ensuremath{{}^{\prime}}并积分为x(\ensuremath{\phi}), 片势V; (ii)求V(\ensuremath{\sigma}); (iii)势降至$V\e^{-3}$的距离xs随eV/T作图, 与线性结果比较. 2006原PDF第2页的指数为$-3$，与2005同题一致。
\end{exercise}
\begin{solution}
令 $w=e\phi/T$, $\lambda=\sqrt{\epsilon_0T/(2ne^2)}$为双物种线性长. Poisson $w''=\sinh w/\lambda^2$. 无穷远w=w\ensuremath{{}^{\prime}}=0, 首积分 $w'^2=4\lambda^{-2}\sinh^2(w/2)$. 对正片x>0取负根, 得
\begin{equation}
 \tanh\frac{w(x)}4=\tanh\frac{w_s}4\e^{-|x|/\lambda},\quad
 \sigma=\frac{4\epsilon_0T}{e\lambda}\sinh\frac{w_s}2.
\end{equation}
故V=(2T/e)arsinh[\ensuremath{\sigma}e\ensuremath{\lambda}/(4\ensuremath{\epsilon}0T)], 且 $x_s=\lambda\ln[\tanh(w_s/4)/\tanh(w_s/(4\e^3))]$. 小ws时xs\ensuremath{\to}3\ensuremath{\lambda}, 大ws时表面势下降更快, xs/\ensuremath{\lambda}\ensuremath{\to}0; 远处尾仍线性指数. 负片用|ws|同理.
\begin{center}\begin{tikzpicture}\begin{axis}[width=.65\textwidth,height=4cm,xlabel={$|eV/T|$},ylabel={$x_s/\lambda$},xmin=0,xmax=20,ymin=0,ymax=3.3]
\addplot[thick,structurecolor,domain=.001:20,samples=100]{ln(tanh(x/4)/tanh(x/(4*exp(3))))};\addplot[dashed,domain=0:20]{3};\end{axis}\end{tikzpicture}\end{center}
\end{solution}
\begin{exercise}\label{ex:06-1-4}
原题4. 半径a无限长冷等离子体圆柱, Te=Ti=0, ne=ni=n0. 电子柱横向位移\ensuremath{\delta}\ensuremath{\ll}a. (i)离子不动求频率; (ii)离子质量M可动重求. 提示:表面电荷产生二维偶极场.
\end{exercise}
\begin{solution}
均匀柱内部单物种场 $\mathbf E=\rho_q\mathbf r/(2\epsilon_0)$. 两柱中心错开\ensuremath{\delta}, 重叠内部净E为 $en_0\boldsymbol\delta/(2\epsilon_0)$, 电子受反向回复, $\ddot\delta=-\omega_{pe}^2\delta/2$. (i)\ensuremath{\omega}=\ensuremath{\omega}pe/\ensuremath{\sqrt2}. (ii)相对位移\ensuremath{\delta}=\ensuremath{\delta}e\ensuremath{-}\ensuremath{\delta}i的约化质量使 $\omega^2=(\omega_{pe}^2+\omega_{pi}^2)/2$, 中心质量自由平移. 1/2是无限圆柱的退极化系数, 不同球形系数为1/3.
\end{solution}
